\documentclass{article}
\usepackage[UTF8]{ctex}
\usepackage{amsmath,amssymb,amsfonts}
\usepackage{enumitem}
\usepackage{geometry}
\geometry{a4paper, margin=1in}
% === 错题本排版优化（不参与题目标识） ===
\setlength{\parindent}{0pt}
\setlength{\parskip}{0.35em plus 0.10em minus 0.08em}
\setlength{\emergencystretch}{3em}
\setlist{itemsep=0.25em, topsep=0.40em, parsep=0pt}
\allowdisplaybreaks[2]
% === 错题本排版优化结束 ===
\begin{document}
\section*{中国石油大学(华东)2025年数学分析真题}
\section*{29．中国石油大学（华东）2025年数学分析试卷}

\subsection{题目1}
1．求极限 \(\lim _{x \rightarrow 0} \cot x\left(\frac{1}{\sin x}-\frac{1}{x}\right)\) ．

\subsection{题目2}
2．求不定积分 \(\int x \arctan x \mathrm{~d} x\) ．

\subsection{题目3}
3．计算第一类曲面积分 \(\iint_{\sum} \frac{1}{z} \mathrm{~d} S\) ，其中积分曲面 \(\sum\) 是上半球面 \(z=\sqrt{4-x^{2}-y^{2}}\) 被平面 \(z=1\) 所截下的顶部部分。

\subsection{题目4}
4．求幂级数 \(\sum_{n=0}^{\infty}(-1)^{n} \frac{x^{2 n+1}}{2 n+1}\) 的收敛域及和函数 \(S(x)\) ，并求 \(\sum_{n=0}^{\infty}(-1)^{n} \frac{1}{2 n+1}\) 的和．

\subsection{题目5}
5．证明 \(\left\{a_{n}\right\}\) 收敛的充要条件是级数 \(\sum_{n=1}^{\infty}\left(a_{n+1}-a_{n}\right)\) 收敛。

\subsection{题目6}
6．已知数列 \(\left\{a_{n}\right\}\) 非负有界，证明： \(\lim _{n \rightarrow \infty} \sqrt[n]{a_{1}^{n}+a_{2}^{n}+\cdots+a_{n}^{n}}=\sup _{n \geq 1}\left\{a_{n}\right\}\) ．

\subsection{题目7}
7．证明：\(\oint_{L} \frac{x \mathrm{~d} y-y \mathrm{~d} x}{3 x^{2}+4 y^{2}}=\frac{\pi}{\sqrt{3}}\) ，其中积分曲线 \(L\) ：椭圆 \(2 x^{2}+3 y^{2}=1\) ，方向为逆时针。

\subsection{题目8}
8．若 \(f(x)\) 在 \([a, b]\) 内有有限导数，且 \(f^{\prime}(a)<0, f^{\prime}(b)>0\) ，则至少存在 \(c \in(a, b)\) ，使得 \(f^{\prime}(c)=0\) ．

\subsection{题目9}
9．计算 \(I=\iiint_{V}\left(x y^{2}+z\right) \mathrm{d} V\) ，其中 \(V\) 是由曲面 \(z=x^{2}+y^{2}\) 与平面 \(z=4\) 围成的区域．

\subsection{题目10}
10．证明：若函数 \(f(x)\) 在 \([a, b]\) 上可积，函数 \(g(x)\) 在 \([a, b]\) 上除一点 \(x_{0}\) 外 \(f(x)=g(x)\) ，则 \(g(x)\) 在\([a, b]\) 上可积，且 \(\int_{a}^{b} f(x) \mathrm{d} x=\int_{a}^{b} g(x) \mathrm{d} x\) ．

\subsection{题目11}
11．分析函数 \(f(x, y)=\left\{\begin{array}{ll}\frac{x^{4}+y^{4}}{x^{2}+y^{2}}, & (x, y) \neq(0,0) ; \\ 0, & (x, y)=(0,0) .\end{array}\right.\) 在 \((0,0)\) 处的可微性．

\subsection{题目12}
12．设函数 \(f(x)\) 在 \([1,3]\) 上连续，在 \((1,3)\) 内可导，且满足 \(f(1)=\int_{2}^{3} x f(x) \mathrm{d} x\) ，证明：至少存在\(\eta \in(1,3)\) ，使得 \(f(\eta)+\eta f^{\prime}(\eta)=0\) ．

\subsection{题目13}
13．证明含参量积分 \(F(y)=\int_{1}^{+\infty} \sin (x y) e^{-x^{2}} \mathrm{~d} x\) 在 \(y \in[1,2]\) 上可导．

\subsection{题目14}
14．计算曲面积分

\[
\iint_{S} \frac{x^{2} \mathrm{~d} y \mathrm{~d} z+y^{2} \mathrm{~d} z \mathrm{~d} x+z^{2} \mathrm{~d} x \mathrm{~d} y}{\sqrt{x^{2}+y^{2}-z+4}} .
\]

其中 \(S\) 是曲面 \(x^{2}+y^{2}=z\) 介于平面 \(z=0\) 及 \(z=1\) 之间的部分，方向取上侧．

\subsection{题目15}
15．用闭区间套定理证明致密性定理，即设数列 \(\left\{x_{n}\right\}\) 有界，则必有收敛子列．
\end{document}
